\cotefile{59}
\scanpage{2}\mspage{1}

\diagram{\begin{tikzcd}[ampersand replacement=\&,row sep=small]
X\arrow[d,"f"']\\ S
\end{tikzcd}}

$f$ propre plat présentation finie \quad $f_*(\underline O_X)\simeq\underline O_S$ univ!

\diagram{\begin{tikzcd}[ampersand replacement=\&]
X\arrow[r,dash]\arrow[d,dash]\&X_P\arrow[d,dash]\\
S\arrow[r,dash]\&P
\end{tikzcd}}

\[
\underline{\mathrm{Pic}}_{X/S}=P
\]

a) Supposons $X/S$ muni d'une section $g$. Alors sur $X\times_S P$ il y a un faisceau inversible canonique $\mathcal L_g$ [trivialisé sur $g(S)\times_S P$]. On en conclut

\[
(Rf_P)_*(\mathcal L_g)\in D_{\mathrm{parf}}(P)
\qquad [f_P:X\times_S P\mathbin{-}P].
\]

Nous nous intéressons en particulier à

% slanted
\margin{on suppose \unsure{th. de finitude|th. du préambule} \unsure{valables}, p. ex. $f$ projectif ou $S$ loc. noeth. etc.}

\[
\det{}^{*}(Rf_P)_*(\mathcal L_g)=M_g
\]

qui est un module inversible [gradué s. il y a lieu] sur $P$.

b) Remplaçons $g$ par une section $g'$. On a alors un isom canonique

\[
\mathcal L_{g'}=\mathcal L_g\otimes_{\underline O_P}{g'}_P^{*}(\mathcal L_g)^{-1}
\]

d'où

\[
(Rf_P)_*(\mathcal L_{g'})\simeq(Rf_P)_*(\mathcal L_g)\,{g'}_P^{*}(\mathcal L_g)^{-1}
\]

et par suite un isom canonique

\[
M_{g'}\simeq M_g\,{g'}_P^{*}(\mathcal L_g)^{-\chi}=M_g\left(N_gN_{g'}^{-1}\right)^{\chi}.\tag*{(*)}
\]

où $\chi$ est la fonction continue $P\longrightarrow\mathbb Z$ donnée par le rang de $Rf_{P*}(\mathcal L_g)$, i.e. la car. d'EP des faisceaux inversibles sur les fibres de $X$ \ldots.., et où $N$ désigne un faisceau inv. \inserted{``de Weil''} quelconque sur $X\times_{\unsure{P|S}}P$, et $N_g$ (resp. $N_{g'}$) les faisceaux inv. sur $P$ qu'il définit via la section $g$ (resp. $g'$). On notera que le facteur de droite

% margin note written sideways
\margin{$N_gN_{g'}^{-1}$ est \uline{algébriquement équivalent} à zero sur $P$ rel/$S$, donc $M_{g'}$ est alg. équiv. à $M_g$ \unsure{rel/$S$|polaris\ldots}.}

c) Soit $\delta$ une section de $\underline{\mathrm{NS}}_{X/S}$, qui définit donc un torseur $P^\delta$ sous $P^{\unsure{\circ|0}}=\underline{\mathrm{Pic}}^\circ_{X/S}$, et considérons la restriction de $M_g$, \inserted{$M_{g'},\ldots$} à $P^\delta$. Nous supposons pour simplifier, par la suite, que $P^\circ$ est un schéma \uline{abélien} sur $S$.
\scanpage{3}\mspage{2}

Notons que $M_g$ et $M_{g'}$ définissent le \={m} élément de $\underline{\mathrm{NS}}_{P^\delta/S}(S)$ en vertu de b), or on a un isom canonique

\[
\underline{\mathrm{NS}}_{P^\delta/S}\simeq\underline{\mathrm{NS}}_{P^{\unsure{\circ|0}}/S}.
\]

Donc on trouve une section

\[
\varphi(\delta)\in\underline{\mathrm{NS}}_{P^\circ/S}(S)
\]

indépendante du choix de la section $\deleted{\unsure{\delta|g}}g$ de $X/S$. Par descente, elle est définie \={m} sans supposer qu'il existe une section. Comme la formation de $\varphi(\delta)$ est compatible (par construction) à l'extension de la base, on trouve en fait un homomorphisme de schémas (pas nécessairement additif)

% slanted
\margin{\uline{détailler}}

\[
\boxed{\varphi:\underline{\mathrm{NS}}_{X/S}\longrightarrow\underline{\mathrm{NS}}_{\deleted{P}B/S}}\qquad(B=P^\circ=\underline{\mathrm{Pic}}^\circ_{X/S}).
\]

d) Revenons à $\delta$\deleted{\illegible{1}}. Notons que $\underline{\mathrm{Pic}}_{P^\delta/S}$ est déduit de $\underline{\mathrm{Pic}}_{\deleted{P}B/S}$ en tordant par \deleted{l'action} le torseur $P^\delta$ sous $P^\circ$, via l'action naturelle de $\deleted{P}B$ sur $\underline{\mathrm{Pic}}_{\deleted{P}B/S}$ déduite de l'action de $\deleted{P}B$ sur lui-= par translation. \deleted{\unsure{Si $\lambda$|Si le}}

Cette action est triviale sur $\underline{\mathrm{Pic}}^\circ_{\deleted{P}B/S}$ et $\underline{\mathrm{NS}}_{B/S}$ [d'où en particulier un isom

\[
\underline{\mathrm{Pic}}^\circ_{P^\delta/S}\simeq\underline{\mathrm{Pic}}^\circ_{B/S},\qquad
\underline{\mathrm{NS}}_{P^\delta/S}\simeq\underline{\mathrm{NS}}_{B/S}\quad].
\]
\scanpage{4}\mspage{3}

Si $\deleted{\unsure{\delta|g}}\lambda$ est une section de $\mathrm{NS}_{\deleted{P}B/S}$ [donc définit une section, notée encore $\lambda$, de $\underline{\mathrm{NS}}_{P^\delta/S}$], alors $\underline{\mathrm{Pic}}^\lambda_{P^\delta/S}$ se déduit de $\underline{\mathrm{Pic}}^\lambda_{B/S}$ en tordant à l'aide de $P^\delta$. Or soit

\[
A=\underline{\mathrm{Pic}}^{\unsure{\circ|0}}_{B/S}\quad(=\underline{\mathrm{Alb}}^\circ_{X/S})
\]

le schéma abélien dual de $B$, \deleted{\unsure{$(=\underline{\mathrm{Alb}}^\circ($}} alors $\lambda$ définit un hom

\[
\widetilde\lambda:B\longrightarrow A
\]

et l'action de $B$ sur $\underline{\mathrm{Pic}}^\lambda_{B/S}$ se déduit de l'action évidente de $A=\underline{\mathrm{Pic}}^\circ_{B/S}$ sur $\underline{\mathrm{Pic}}^\lambda_{B/S}$ (qui est un $A$-torseur) via l'hom $\widetilde\lambda$. Donc on trouve

\[
\underline{\mathrm{Pic}}^\lambda_{P^\delta/S}\simeq\underline{\mathrm{Pic}}^\lambda_{B/S}
\overset{A}{\times}
\underbrace{\left(P^\delta\overset{B}{\times} A\right)}_{
\substack{A\text{-torseur déduit}\\\text{du }B\text{-torseur}\\ P^\delta\text{ grâce à l'hom.}\\ \widetilde\lambda:B\to A}}
\]

\deleted{\unsure{No|Alor|Alo}.} Reprenant les notations de c), nous allons appliquer ceci au cas où l'on a

\[
\lambda=\varphi(\delta),
\]

de sorte que, pour \deleted{\unsure{la|tt}} section $g$ de $X/S$, la restriction $L_g\mid P^\delta$ définit une section de $\underline{\mathrm{Pic}}^\lambda_{P^\delta/S}$, d'où une section
\scanpage{5}\mspage{4}

\[
\ell_\delta(g)\in\underline{\mathrm{Pic}}^{\deleted{\unsure{\lambda}}\varphi(\delta)}_{B/S}
\overset{A}{\times}(P^\delta\overset{B}{\times}(A,\widetilde\lambda))
\]

La formation de cette dernière commute évidemment au changement de base, donc on a défini un morphisme de schémas (ou plutôt, de foncteurs\ldots)

\[
X\xrightarrow{\ell_\delta}\underline{\mathrm{Pic}}_{B/S}^{\deleted{\unsure{\lambda}}\varphi(\delta)}
\overset{A}{\times}(P^\delta\overset{B}{\times}(A,\widetilde{\deleted{\unsure{\lambda}}\varphi(\delta)}))\ , 
\]

\margin{\diagram{\begin{tikzcd}[ampersand replacement=\&,column sep=small,row sep=tiny]
A\arrow[d,equal]\&B\arrow[l,"\mid"]\arrow[d,equal]\\
\mathrm{Alb}^{\unsure{\circ|0}}\&\mathrm{Pic}^\circ
\end{tikzcd}}}

Comme le deuxième mb est un torseur sous le schéma abélien $A$, on trouve \inserted{une factorisation de $\ell_\delta$ par} un homomorphisme

\diagram{\begin{tikzcd}[ampersand replacement=\&,row sep=small]
\underline{\mathrm{Alb}}^1_{X/S}\deleted{(X/S)}\arrow[r]\&
\underline{\mathrm{Pic}}_{B/S}^{\deleted{\unsure{\lambda}}\varphi(\delta)}\overset A\times (P^\delta\overset B\times(A,\widetilde{\deleted{\unsure{\lambda}}\varphi(\delta)}))\\
A\arrow[r]\& A
\end{tikzcd}}

qui est \deleted{un homomorphisme} \inserted{compatible avec un} certain homomorphisme $u_\delta:A\longrightarrow A$, sur les Groupes structuraux $A=\underline{\mathrm{Alb}}^\circ_{X/S}\deleted{\unsure{(X/S)}}$. Pour voir quel il est, on applique la formule \inserted{$(*)$} de b), on trouve $u_\delta=-\chi(\delta)\,\mathrm{id}_A$, où $\chi=\chi(\delta)$ est la fonction $S\longrightarrow\mathbb Z$ dont l'image inverse sur $P^\delta$ est $\chi\mid P^\delta$ ($\chi$ comme dans b)). On trouve donc un \deleted{$A$} homomorphisme de $A$-torseurs

% slanted
\margin{vérifier si le signe est correct !}

\[
\underline{\mathrm{Alb}}^{-\chi(\delta)}_{X/S}\longrightarrow
\mathrm{Pic}^{\varphi(\delta)}_{B/S}\overset{A}{\times}
(P^\delta\overset B\times(A,\widetilde{\varphi(\delta)}))
\]

i.e. une \uline{section canonique}

\[
\boxed{L(\delta)\in\Gamma\bigl(S,\underline{\mathrm{Alb}}^{\chi(\delta)}_{X/S}
\overset A\times\mathrm{Pic}_{B/S}^{\varphi(\delta)}\times
(P^\delta_{X/S}\overset B\times(A,\widetilde{\varphi(\delta)}))\bigr)}
\]
\scanpage{6}\mspage{5}

À fortiori, \uline{le $A$-torseur}

\[
\underline{\mathrm{Alb}}^{\chi(\delta)}_{X/S}\overset A\times\underline{\mathrm{Pic}}^{\varphi(\delta)}_{B/S}\overset A\times(P^\delta\overset B\times(A,\widetilde{\varphi(\delta)}))
\]

\uline{est trivial} (en fait, avec une trivialisation canonique).

e) Pour exploiter ces considérations, il \deleted{\unsure{faudra,|faudrait}} encore, pour tout $\delta\in\mathrm{NS}(X/S)$, ``calculer'' $\varphi(\delta)$ (et $\chi(\delta)$ \ldots). Pour ceci, on \deleted{peut supposer} est ramené au cas $S=\mathrm{Spec}\,k$, \unsure{($k$|$k$|et $k$} un corps, alg. clos si on y tient\deleted{\unsure{)}}. Choisissons alors $p\in P^\delta(k)$, qui correspond à un faisceau inversible $(\mathcal L_g)_p=L$ sur $X$, dans la classe $\delta$.

\margin{$X\times P^\delta$}

\unsure{Le choix|La donnée} de $\deleted{\unsure{g}}p$ définit un isom

\[
P^\delta\simeq B
\]

et moyennant celui-ci, le faisceau $\mathcal L_g\mid X\timesP^\delta$ \deleted{$\unsure{X\times P^\delta}$} est transformé en \deleted{\unsure{$\mathcal L_g$ de $X\times$}} le faisceau $\mathcal L^{\unsure{0|\circ}}\deleted{\unsure{\otimes_X|\otimes}}\underline{\mathrm{pr}}_1^*(L_{\unsure{\delta|{}}})$, où $\mathcal L^\circ$ est le faisceau de Weil sur $X\times B$ [relatif au choix du point $g\in X(k)$) et où $L$ est \deleted{\unsure{de}}le faisceau inversible sur $X$ \deleted{défini} de classe $\delta$ déjà envisagé. Il s'agit donc de calculer

\[
\det{}^{*}(R\mathrm{pr}_{2*})(\mathcal L\otimes\mathrm{pr}_1^*(L))
\]
\scanpage{7}\mspage{6}

\deleted{Si on fait} mod équivalence algébrique. \deleted{Comme $\mathrm{NS}(B/S)$ est} \deleted{\unsure{Évidemt}} \deleted{sans torsion, il suffit de faire} \deleted{(cela suffit : \unsure{Not}} le calcul dans $\mathrm{NS}(B/S)\otimes_{\mathbb Z}\mathbb Q$, \deleted{et à fortiori dans \unsure{et}} on trouve alors le terme de degré 1 dans $\mathrm{ch}_B(R\mathrm{pr}_{2*}(\mathcal L\otimes\mathrm{pr}_1^*(L)))$. Or le $\mathrm{ch}_B$ se calcule par RR, du moins si $X$ est loc une intersection complète, on trouve

\[
\begin{aligned}
\varphi(\delta)&=\text{terme en degré 1 dans}\\
&\mathrm{pr}_{2*}\left(\exp(D+\mathrm{pr}_1^*\Delta)\,\mathrm{pr}_1^*(\mathrm{Todd}_{X/k})\right)
=\mathrm{pr}_{2*}^{\deleted{\unsure{*}}}\bigl(\deleted{\unsure{(|\mathrm{pr}}}\inserted{\exp D}\,\mathrm{pr}_1^*(T_{X/k}\exp\Delta)\bigr)\\
&=[\exp D](T_{X/k}\exp\Delta)
\end{aligned}
\]

\[
\text{où}\quad\left\{\begin{aligned}
D&=\mathrm{cl}\,\mathcal L\in\mathrm{Gr}^1(X\times B)\\
\Delta&=\mathrm{cl}\,\Delta\in\mathrm{Gr}^1(X)\\
\mathrm{Todd}_{X/k}&\in\mathrm{Gr}^*(X)
\end{aligned}\right.
\]

Nous allons faire le calcul en supposant $X$ lisse de dimension $\unsure{n|h}$, et en travaillant en cohomologie $l$-adique ($l.1_k\ne0$). On choisit, pour simplifier, un isomorphisme $T_l(\unsure{\bar k^*|k^*})\simeq\mathbb Z_l$, pour éviter les \unsure{twists|Twists|twist} de Tate. On a

\[
D\in H^2(X\times B)=H^2(X)+H^1(X)\otimes H^1(B)+H^2(B)
\]

et en fait $\deleted{\unsure{D}}$ on a

\[
D\in H^1(X)\otimes H^1(B)
\]

(et $D$ définit un isomorphisme entre $H^1(B)$ et $H^1(X)^\vee$\unsure{)|,|),}
\scanpage{8}\mspage{7}

On peut écrire

\[
D=\sum \unsure{e_i|e^i}\otimes e'_i
\]

où $(e_i)$ est une base de $H^1(X)$, et $e'_i$ est la base duale de $H^1(B)$. Posant

\[
\mathrm{Todd}_{X/k}\exp\Delta=\sum_{i=0}^{n}a^i\qquad a^i\in H^{2i}(X)
\]

on trouve

\[
\exp D\,\mathrm{pr}_1^*(\mathrm{Todd}_{X/k}\exp\Delta)=\sum_{i,j}\mathrm{pr}_1^*(a^i)\,D^j/j!
\]

le terme en $i,j$ écrit étant de degré cohomologique partiel en le facteur $X$ égal à $2i+j$. Il donne une contribution nulle dans $\mathrm{pr}_{2*}$ sauf si $2i+j=2n$ i.e. $j=2(n-i)$, \deleted{donc il faut} \inserted{auquel cas on trouve} un $\mathrm{pr}_{2*}$ de degré $j$, qui n'est intéressant que si $j=2$ i.e. $n-i=1$. \deleted{prendre} Donc on trouve

\[
\deleted{\unsure{\varphi(\delta)=|\sum}\tfrac12\mathrm{pr}_1^*(a_{n-1})D^2\unsure{\in H^2(B)}.}
\]

\[
\varphi(\delta)\deleted{\unsure{\in}}=\tfrac12\mathrm{pr}_{2*}\bigl(\mathrm{pr}_1^*(a_{n-1})D^2\bigr)
\]

\uline{Ex. 1}

Prenons le cas $n=1$. Alors $\delta$ est caractérisé par son degré, soit \deleted{\unsure{$\delta$}} $d$. On trouve

\[
\deleted{\mathrm{Todd}_{X/k}\exp\Delta=(1\unsure{-|+}\tfrac K2)(1+\Delta)=1+\underbrace{(\Delta-\tfrac K2)}_{\substack{\text{degré}\\d-g+1}}\unsure{\nearrow|\diagup}}
\]

\[
\varphi(\delta)=\tfrac12\mathrm{pr}_{2*}(D^2)
\]

\deleted{et il faut prendre $\mathrm{pr}_{2*}$ de}

\[
\deleted{\tfrac12 D^2+\mathrm{pr}_1^*(\Delta-\tfrac K2)}
\]

Or \deleted{\unsure{nous}} on a, si $(e_1,f_1,\unsure{e_2,f_2|e,f},\ldots,e_g,f_g)$ est une base symplectique de $H^1(X)$ [$e_i f_i=-f_i e_i=\eta$, base can. de $H^2(X)$),

% slanted
\margin{On identifie $H^1(X)$ à son dual grâce à la forme \unsure{intersection|d'intersection}.}

\[
D=\sum(e_i\otimes f_i+f_i\otimes e_i)
\]
\scanpage{9}\mspage{8.}

donc

\[
\begin{aligned}
D^2&=\sum_{i,j}\underbrace{(e_i\otimes f_i+f_i\otimes e_i)(e_{\unsure{j|g}}\otimes f_j+f_j\otimes e_j)}_{-[e_i f_j\otimes f_i e_j+f_i e_j\otimes e_i f_j]}\\
&=-\sum_i(\eta\otimes f_i e_i-\eta\otimes e_i f_i)=2\eta\sum e_i f_i
\end{aligned}
\]

\[
\text{donc}\qquad\tfrac12D^2=\eta\otimes\sum e_i f_i
\]

et par suite

\[
\varphi(\delta)=\sum e_i\wedge f_i\in H^2(B)
\]

\uline{C'est la classe de polarisation canonique} de la \uline{jacobienne $B$}. D'ailleurs $\chi(\delta)=1-g+d$

% slanted
\margin{vérifier que le signe est correct \unsure{ici}}

\deleted{\unsure{R|et}} Utilisons la polarisation canonique \inserted{par le diviseur $\Theta$} de $B$, définissant un isom can

\[
c:B\longrightarrow A
\]

et transformons les $A$-torseurs de la fin de d) par $c^{-1}$. \uline{On sait que} $\underline{\mathrm{Alb}}^1_{X/S}$ \uline{est transformé} en $\underline{\mathrm{Pic}}^{-1}_{X/S}=P^{-1}$ (\uline{attention au signe} !) donc $\underline{\mathrm{Alb}}^{\chi(\delta)}_{X/S}$ devient $P_{\deleted{\unsure{X/S}}}^{-\chi(\delta)}=P^{g-1-d}$, d'autre part $\underline{\mathrm{Pic}}^{\varphi(\delta)=\text{pol. can. de } A}_{B/S}$ \deleted{canoniquement} \inserted{devient} $\deleted{\unsure{\mathrm{Pic}^{g-1}|\underline{\mathrm{Pic}}^1}}\ P^{g-1}$ (isom. canonique respectant les div. $\Theta$), enfin $P^\delta\overset B\times(A,\underset{\substack{\shortparallel\\c}}{\widetilde{\varphi(\delta)}})$ devient $P^\delta=P^d$. Donc le torseur envisagé, produit \inserted{contracté} des trois précédents, est $P^{2g-2}$. On trouve une section canonique de ce dernier : ce n'est pas très excitant, car on s'attend à trouver la section canonique $K$ (à prouver !).
\scanpage{10}\mspage{9}

\uline{NB} Pour bien faire, il faudrait prouver que sur le site modulaire $S$ des schémas en courbes \inserted{relatives propres lisses connexes} de genre $g$\deleted{$\unsure{(g\geq1)|\geq1}$} \inserted{($g\geq1$ ; on \deleted{\unsure{distingue}}\inserted{\unsure{fixe|prend|choisit}} une section marquée pour $g=1$)} \deleted{\unsure{le|tel que}} si $\mathcal X$ est la courbe de genre $g\deleted{\unsure{*}}$ universelle, $P=\underline{\mathrm{Pic}}_{X/S}$, alors le groupe des sections de $P$ sur $S$ est engendré par $K$ (\deleted{donc} \inserted{en particulier, que} le groupe des sections de $P^{\unsure{\circ|0}}$ sur $S$ est égal à zéro). Il suffirait de déterminer le groupe en question en car. nulle.

\deleted{\unsure{f)}} \inserted{f) \uline{Ex. 2}} \uline{Cas} $X$ \uline{un schéma abélien}. Donc $X=A$. On sait que \inserted{sur fibres géométriques} $\exp D$ définit $\deleted{\unsure{A\to A}}\inserted{\Phi}:\mathrm{Chow}(A)_{\mathbb Q}\xrightarrow{\sim}\mathrm{Chow}(B)_{\mathbb Q}$ qui \deleted{\unsure{en}} transforme élément de degré $i$ en élément de degré complémentaire $n-i$, \inserted{et} qui transforme produit ordinaire en produit de convolution, et inversement. D'ailleurs $T_{X/k}=1$, donc on trouve

\[
\varphi(\delta)=\inserted{\text{composante de degré 1 de }}\sum\tfrac1{i!}\underbrace{\Phi(\Delta)}_{\substack{\text{classe d'une}\\\text{courbe}}}^{*i}
=\frac1{(n-1)!}\Phi(\Delta)^{*(n-1)}.
\]

\deleted{\unsure{donc}} On trouve ainsi un hom (\inserted{polynomial} de degré $n-1$)

\[
\varphi:\underline{\mathrm{NS}}_{\deleted{\unsure{X/S}}A/S}\longrightarrow\underline{\mathrm{NS}}_{B/S}.
\]

On vérifie que, si $\delta$ est une section de $\underline{\mathrm{NS}}_{A/S}$, définissant $\widetilde\delta:A\longrightarrow B$ d'où $\varphi(\delta)$ section de $\underline{\mathrm{NS}}_{\unsure{B|P}/S}$, d'où

\[
\widetilde{\varphi(\delta)}:B\longrightarrow A
\]

\margin{faire calcul direct, en faisant \deleted{\unsure{attention}} attention aux signes !}

\[
\text{on a}\qquad\boxed{\widetilde{\varphi(\delta)}\,\widetilde\delta=-\chi(\delta)\,\mathrm{id}_A,\qquad\widetilde\delta\,\widetilde{\varphi(\delta)}=-\chi(\delta)\,\mathrm{id}_B}
\]
\scanpage{11}\mspage{10}

Par suite, si $\delta$ est non dégénéré i.e. $\chi(\delta)\ne0$, alors $\varphi(\delta)$ est non dégénérée, \deleted{et si $\delta$ est une \unsure{polar|polarisation}} \inserted{et si $\delta$ est une} polarisation, $-\varphi(\delta)\inserted{=\delta'}$ est une polarisation. On \inserted{déduit} d'ailleurs de la relation plus haut, \deleted{\unsure{que|puisque}} $\deg\widetilde\delta=\chi(\delta)^2$, $\deg\widetilde{\delta'}=\chi(\delta')^2$, $\deg\widetilde\delta\,\widetilde{\delta'}=\chi(\delta)^2\chi(\delta')^2$,

\[
\delta'=-\varphi(\delta)\qquad\underset{\chi(\delta)^{2\unsure{n|g}}}{\shortparallel}
\]

% slanted
\margin{$\delta'=-\varphi(\delta)$ à priori, on a égalité des carrés ; mais les deux côtés sont des fonctions polynomiales, sur $\mathrm{NS}$ positives si $\delta$ est une polarisation}

\[
\chi(\deleted{\unsure{\varphi(}}\delta'\deleted{\unsure{)}})=\chi(\delta)^{n-1}
\]

si $\deleted{\unsure{\varphi(}}\delta)$ est non dégénéré, \deleted{alors $\delta$ \unsure{l'est}} $\chi(\delta)$ l'est.

Notons aussi que si $\delta'=-\varphi(\delta)$, $\delta''=-\varphi(\delta')$, on a

\[
\widetilde{\delta'}=\delta^{-1}\chi(\delta),\qquad
\deleted{\varphi(\varphi(\delta))=}\delta''{}^{-1}=\delta'{}^{-1}\chi(\delta')=\delta\chi(\delta)^{-1}\chi(\delta')=\delta\chi(\delta)^{n-2}
\]

\[
\text{donc}\qquad\varphi(\varphi(\delta))=\delta\chi(\delta)^{n-2}
\]

Si \deleted{\unsure{$g$}} $n=1$, i.e. $X$ une courbe elliptique, on trouve \deleted{$\varphi(\delta)={}$ polarisation canonique,} $\delta\chi(\delta)$ i.e. $\varphi(\varphi(\delta))$ est la polarisation canonique ; \deleted{[NB $\varphi($} mieux, pour tt $\delta$, $\varphi(\delta)$ est la polarisation canonique. Si $n\geq2$, la formule $\varphi(\varphi(\delta))=\delta\chi(\delta)^{n-2}$ inspire confiance directement\ldots.

\uline{NB} Il faudrait prouver que sur le \deleted{\unsure{schéma|schéma de}} \deleted{\unsure{topos}} \inserted{site} modulaire des schémas abéliens polarisés $(A,\delta)$ \inserted{en car. 0}, avec degré de polarisation donné $\chi(\delta)$, la polarisation $\varphi(\delta)=\delta'$ \deleted{\unsure{ind}} sur $B$ n'est plus divisible, \deleted{\unsure{ce que l'hom|i.e. que \illegible{1}}} ce qui montrerait que $\varphi(\delta)$ est ce qu'on peut faire de mieux dans ce genre comme ``inversion'' d'une polarisation. \deleted{\unsure{donc|On}} \inserted{En fait,} on s'attend à ce que $\varphi(\delta)$ engendre le groupe \deleted{par \unsure{regarder au pt de vue}} des sections de $\underline{\mathrm{NS}}_{B/S}$ sur une quelconque composante connexe du \deleted{\unsure{site modulaire}} \deleted{topos} \inserted{site}, modulaire\ldots
\scanpage{12}\mspage{11}

L'isomorphisme de la fin de d) nous fournit une section canonique de

\[
\underline{\mathrm{Pic}}_{B/S}^{-\delta'}\overset A\times
(\underline{\mathrm{Pic}}^\delta_{A/S}\overset B\times(\deleted{\unsure{B}}A,-\widetilde{\delta'}))
\]

où $\delta'=-\varphi(\delta)$ comme dessus. En d'autres termes, on trouve un isomorphisme canonique

\[
\boxed{\underline{\mathrm{Pic}}^{\delta'}_{B/S}\simeq\underline{\mathrm{Pic}}^{\deleted{\unsure{=|-}}-\delta}_{A/S}\overset B\times(A,\widetilde{\delta'})}\tag*{(*)}
\]

\margin{$ {}_2 A$}

% slanted
\margin{\uline{NB} Comme $\underline{\mathrm{Pic}}^{-\delta}_{A/S}\simeq\underline{\mathrm{Pic}}^\delta_{A/S}$ (isom canonique) on peut supprimer si on veut les signes dans notre formule\ldots{} Prouver qu'après passage aux carrés des torseurs, on trouve un isom qui est compatible avec les sections canoniques, donc que $(*)$ provient d'un isom entre les torseurs \unsure{correspondants} au Groupe $ {}_2 A$\ldots. \uline{à vérifier}}

\uline{Exemple} Supposons que $\delta$ soit de degré 1, \deleted{\unsure{$\delta'$ est u}} $\widetilde\delta:A\xrightarrow{\sim}B$, et $\widetilde{\delta'}:B\simeq A$ est l'inverse de $\widetilde\delta$. Alors $\delta'$ \deleted{et \unsure{$\delta'$ c'est la}} \inserted{est la} polarisation déduite de $\delta$ par transport de structure à l'aide de $\widetilde\delta$. On trouve un isom., \unsure{grâce|via} à $\widetilde\delta$, $\underline{\mathrm{Pic}}^{\delta'}_{B/S}\simeq\underline{\mathrm{Pic}}^\delta_{A/S}$ (compatible \deleted{\unsure{avec $\underline{\mathrm{Pic}}$}} avec $\deleted{\unsure{B}}A\xrightarrow[\sim]{\widetilde\delta}B$) et l'isom considéré peut s'interpréter comme un isom \inserted{de $B$-torseurs}

\[
\underline{\mathrm{Pic}}^\delta_{A/S}\simeq\underline{\mathrm{Pic}}^{-\delta}_{A/S}
\]

\deleted{C'est là l'isomorphisme canonique bien connu, défini par la condition de \unsure{transformer} le diviseur $\Theta^{\unsure{-|{}}}$ en $\Theta^{-}$.} \deleted{\unsure{sur $A\times B$|$A\times B$}} \deleted{\unsure{Je crois que c'est la propriété}\illegible{2}\ldots}

\uline{Remarque} La section du torseur de la fin de d) est stable par autom. de $(X/S,\delta)$ ; dans le cas actuel, on peut prendre les autom par translations dans $X=A$. On trouve que la section par $a$ opère sur

\[
\underline{\mathrm{Alb}}_{X/S}^{\chi(\delta)}\times\underline{\mathrm{Pic}}^{\varphi(\delta)}_{B/S}
\overset A\times (P^\delta\overset B\times(A,\widetilde{\varphi(\delta)}))
\simeq\underline{\mathrm{Pic}}^{\varphi(\delta)}_{B/S}\overset A\times
(P^\delta\overset B\times(A,\widetilde{\varphi(\delta)}))
\]

% slanted
\margin{Attention au signe !?!}

par translation par $\chi(\delta)a+\widetilde{\varphi(\delta)}\,\widetilde\delta(a)\deleted{\unsure{a}}$ ; donc on a $[\chi(\delta)\mathrm{id}_A+\widetilde{\varphi(\delta)}\,\widetilde\delta]a=0$ pour tt $a$, d'où \deleted{\unsure{résulte}} $\widetilde{\varphi(\delta)}\,\widetilde\delta=-\chi(\delta)\mathrm{id}_A$
\scanpage{13}\mspage{12}

\uline{Remarque} Dans l'isom can \unsure{$(*)$}, remplaçons $\delta$ par $\lambda\delta$ ($\lambda\in\mathbb Z$). Les deux mbs sont remplacés respectivement par $\underline{\mathrm{Pic}}^{\lambda^{n-1}\delta'}_{B/S}\simeq(\underline{\mathrm{Pic}}^{\delta'}_{B/S})^{(\lambda^{n-1})}$, et par $\underline{\mathrm{Pic}}^{-\lambda\delta}_{A/S}\overset B\times(A,\lambda^{n-1}\widetilde{\delta'})\simeq(\underline{\mathrm{Pic}}^{-\delta}_{A/S}\overset B\times(A,\widetilde{\delta'}))^{(\lambda^n)}$. Donc on trouve, grâce à l'égalité des deux mbs pour $\lambda\delta$, que

\[
(\underline{\mathrm{Pic}}^{\delta'}_{B/S})^{\otimes(\lambda^n-\lambda^{n-1})}\simeq\text{torseur trivial}\quad\text{(canoniqt)}
\]

Faisant $\lambda=-1$, on trouve

\[
(\underline{\mathrm{Pic}}^{\delta'}_{B/S})^{\otimes2}\simeq\text{torseur trivial}\quad\text{(canoniqt)}
\]

% slanted
\margin{Vérifier que la trivialisation est la trivialisation canonique définie pour \uline{toute} section $\delta'$ de $\underline{\mathrm{NS}}_{B/S}$}

\uline{NB} 2 est le ppcm des entiers de la forme $\lambda^n-\lambda^{n-1}$].

\deleted{Lorsque} Comme $(\delta')'=\chi(\delta)^{n-2}\delta$, on en conclut, en échangeant les rôles de $A$ et $B$

\[
(\underline{\mathrm{Pic}}^\delta_{\deleted{\unsure{B}}A/S})^{\otimes(2\chi(\delta)^{n-2})}\simeq\text{torseur trivial}\quad\text{(canon.!)}
\]

% slanted
\margin{Si $n\geq2$ (tout est trivial si $n\leq1$)}

Lorsque $\chi(\delta)$ est pair, on a même

\[
(\underline{\mathrm{Pic}}^\delta_{A/S})^{\otimes(\chi(\delta)^{n-2})}\simeq\text{torseur trivial}\quad\text{(canoniqt)}
\]

pour $\chi(\delta)$ \uline{pair}, comme on voit par les réflexions qui suivent\ldots

% slanted
\margin{Si $n\geq3$, vérifier que cette section est encore la section canonique de $(\underline{\mathrm{Pic}}_{A/S}^{\delta\chi(\delta)^{n-2}/2})^{\otimes2}$. Si $n\geq2$, $\chi(\delta)$ pair implique que $\delta$ divisible par 2, et il faudrait prouver encore que la trivialisation est la trivialisation canonique.}

\deleted{Lorsque $\chi(\delta)=1$, et que $\delta$ est une polarisation, il faudrait prouver que la\deleted{\unsure{ite}} section de $(\underline{\mathrm{Pic}}^\delta_{A/S})^{\otimes2}$ obtenue correspond à l'isom $\underline{\mathrm{Pic}}^\delta\simeq\underline{\mathrm{Pic}}^{-\delta}$ transformant $\Theta$ en $\Theta^{-}$\ldots\unsure{.}}
\scanpage{14}\mspage{13}

\deleted{g) Généralisation des constructions précédentes\ldots..} \inserted{\uline{\unsure{Remarque|Remarques}}} Soit $\delta$ \deleted{polarisation de} \inserted{section de $\underline{\mathrm{NS}}_{A/S}$} \deleted{$A$}, d'où $\delta'$ \deleted{pol. de} \inserted{pour} $B$, $\delta''$ \deleted{polarisation de} \inserted{pour} $A$. On a successivement

\[
\underline{\mathrm{Pic}}^{\delta'}_{B/S}\simeq\underline{\mathrm{Pic}}^\delta_{A/S}\overset B\times(A,\widetilde{\delta'})
\]

\[
\underline{\mathrm{Pic}}^{\delta''}_{A/S}\simeq\underline{\mathrm{Pic}}^{\delta'}_{B/S}\overset A\times(B,\widetilde{\delta''})\simeq\mathrm{Pic}^\delta_{A/S}\overset B\times(B,\widetilde{\delta''}\,\widetilde{\delta'})\tag*{(*)}
\]

\[
\text{Or on a}\quad\left\{\begin{aligned}
\delta''&=\delta\chi(\delta)^{n-2}\\
\widetilde{\delta''}\,\widetilde{\delta'}&=\chi(\delta')\,\mathrm{id}_B=\chi(\delta)^{n-1}\,\mathrm{id}_B
\end{aligned}\right.
\]

donc les \unsure{\={m}bs|mbs} extrêmes de $(*)$ donnent

\[
\underline{\mathrm{Pic}}_{B/S}^{\delta\chi(\delta)^{n-2}}\simeq
\underline{\mathrm{Pic}}_{B/S}^{\delta\chi(\delta)^{n-1}}\qquad\text{d'où}
\]

\[
(\underline{\mathrm{Pic}}^\delta_{B/S})^{\otimes[\chi(\delta)^{n-1}-\chi(\delta)^{n-2}]}=\text{torseur trivial}\quad\text{(isom canonique}
\]

Remplaçons $\delta$ par $\lambda\delta$ ($\lambda\in\mathbb Z$) on trouve plus généralement

\[
(\underline{\mathrm{Pic}}^\delta_{B/S})^{\deleted{\unsure{\lambda^n|\otimes}}\chi(\delta)^{n-2}(\lambda^{n}\chi(\delta)\deleted{\unsure{-\lambda}}-\lambda^{n-1})}\simeq\text{torseur trivial}
\]

\[
\deleted{\lambda^{n-1}(\lambda\chi(\unsure{\delta})\unsure{-1)}\ \unsure{\lambda^n}\chi(\delta)^{n-2}[\lambda^n\chi(\delta)-\lambda^{n-1}]}
\]

\deleted{Soit $\alpha$ le \deleted{pgcd} ppcm des $\lambda^n\chi(\delta)-\lambda^{n-1}$ ($\lambda\in\mathbb Z$).} Pour qu'un nb premier $p$ divise $\alpha$ sss on a

\[
\lambda^n\chi(\delta)\deleted{\unsure{\Delta}}-\lambda^{n-1}\equiv0\quad(p)\quad\text{pour }\forall\ \lambda\in\mathbb Z
\]

Faisant $\lambda=\pm1$ et opérant un retranchement on trouve

\[
(\underline{\mathrm{Pic}}^\delta_{B/S})^{\otimes(2\chi(\delta)^{n-2})}\simeq\text{torseur trivial}
\]

\[
\text{et \unsure{si|de} même}\quad(\underline{\mathrm{Pic}}^\delta_{B/S})^{\otimes\chi(\delta)^{n-2}}\simeq\text{torseur trivial si }\chi(\delta)\text{ pair}
\]

[car alors $\chi(\delta)-1$ est impair, et on utilise que

\[
(\underline{\mathrm{Pic}}^\delta_{B/S})^{(\chi(\delta)-1)\chi(\delta)^{n-1}}\simeq\text{torseur trivial}.
\]

\uline{Exemple} Si $n=2$, on trouve : $(\underline{\mathrm{Pic}}^\delta_{B/S})^{\otimes2}$ trivial, et $\mathrm{Pic}^\delta$ trivial si $\chi(\delta)$ pair\ldots
